% --- ARCHIVO: Anexos_Conceptos/anexo_cap5.tex ---

\chapter{Conceptos técnicos y regulatorios del Capítulo 5}
\label{anexo:conceptos_cap5}

\section{Criterio N-1}
\label{anexo:criterio_n1}
El \textit{Criterio N-1} es la norma de seguridad fundamental en la operación y planificación de sistemas eléctricos de potencia. Establece que el sistema debe ser capaz de mantener los parámetros de tensión y frecuencia dentro de los límites operativos normativos tras la pérdida contingente de cualquier elemento único (un generador, una línea de transporte o un transformador), sin provocar cortes de suministro en cascada ni daños en los equipos.

\section{Régimen de Renovables, Cogeneración y Residuos (RCR)}
\label{anexo:rcr}
El Régimen de Renovables, Cogeneración y Residuos (RCR) es el marco regulatorio del sistema eléctrico español que agrupa a las instalaciones de producción de energía eléctrica a partir de fuentes descarbonizadas. Históricamente, las obligaciones técnicas operativas exigidas a este régimen (como el control de tensión o la aportación de servicios de ajuste) han diferido sustancialmente de las exigidas al parque generador síncrono convencional.

\section{Procedimiento de Operación 7.4 (P.O. 7.4)}
\label{anexo:po74}
El P.O. 7.4 es la normativa técnica del sistema eléctrico español que regula el servicio de ajuste de control de tensión en la red de transporte. Define las obligaciones de los generadores para absorber o inyectar potencia reactiva ($Q$) en función de las consignas enviadas por el Operador del Sistema, con el objetivo de mantener los perfiles de tensión ($V$) dentro de los rangos de seguridad estipulados en el P.O. 1.1.

\section{Compensador Síncrono Estático (STATCOM)}
\label{anexo:statcom}
Un STATCOM (\textit{Static Synchronous Compensator}) es un dispositivo de compensación activa de potencia reactiva basado en electrónica de potencia (inversores VSC). A diferencia de las reactancias o los bancos de condensadores mecánicos de conmutación discreta, un STATCOM inyecta o absorbe potencia reactiva de forma continua, dinámica y casi instantánea, proporcionando un control de tensión muy superior frente a transitorios rápidos.

\section{Network Code on Requirements for Generators (NC RfG)}
\label{anexo:nc_rfg}
El \textit{Network Code on Requirements for Generators} es el código de red europeo establecido por ENTSO-E que armoniza los requisitos técnicos obligatorios que deben cumplir las instalaciones de generación para conectarse a la red. Su objetivo es garantizar que todos los generadores, independientemente de su tecnología síncrona o asíncrona, posean las capacidades de resiliencia necesarias para mantener la estabilidad del área síncrona continental.

\section{Inercia sintética}
\label{anexo:inercia_sintetica}
La \textit{inercia sintética} (o inercia virtual) es la capacidad de una instalación de generación basada en inversores (IBR) para emular la respuesta inercial natural de una máquina síncrona tradicional. Mediante algoritmos de control avanzados, el inversor inyecta o absorbe potencia activa de forma ultrarrápida proporcional a la derivada temporal de la frecuencia (RoCoF), mitigando las variaciones bruscas de frecuencia ante desequilibrios de potencia.

\section{Control Grid-forming frente a Grid-following}
\label{anexo:grid_forming}
Este concepto define el paradigma de control de los inversores. Un inversor \textit{grid-following} (seguidor de red) se sincroniza pasivamente con la tensión y frecuencia preexistentes, dependiendo de la red externa para operar. Por el contrario, un inversor \textit{grid-forming} (formador de red) actúa como una fuente de tensión ideal tras una impedancia: establece activamente su propia onda de tensión y frecuencia, permitiendo sostener la red de forma autónoma, aportar robustez (potencia de cortocircuito) y operar incluso con un 100~\% de penetración de electrónica de potencia o en modo isla.

\section{Protecciones de pérdida de sincronismo (OST)}
\label{anexo:ost}
Los relés de pérdida de sincronismo (OST, \textit{Out-of-Step Tripping}) son esquemas de protección sistémica diseñados para detectar divergencias angulares severas entre áreas interconectadas (deslizamiento de polos). Cuando la diferencia de fase angular excede los límites de estabilidad electromecánica y la transmisión de potencia se vuelve insostenible, los relés abren automáticamente las líneas de interconexión para evitar daños estructurales y confinar el colapso, impidiendo su propagación a los sistemas vecinos.